\part{Unresolved comparisons and current gates}
\section{Statements not asserted as proved}

There is no unchecked theorem silently used in the preceding part.
The implications with hypotheses, including B8 and C15's universality
reduction, are proved implications; their hypotheses are not supplied
by this run. The following are the remaining identifications.

\paragraph{G1: the exact iterated character-independent functional.}
The common oscillator kernel, A1's normalized zero orbit, B11's
restriction seesaw, A5's finite-frame identity and A7's actual single
theta lift have been established. B17 now makes the actual wedge
a literal double lift $V\to P\to V$ by the genuine adjoint
kernel and a common positive inverse $Q^{-1}$.
It specializes the analytic last-step Bessel seesaw without a
new occurrence hypothesis. The operator is common to every
character for fixed spaces, kernel and levels.
The remaining \texttt{MainX} question is its stronger arithmetic
uniform relation-preservation statement. To apply it, one must
construct its actual common geometry,
tuples, archimedean maps, definite class kernels and arithmetic links,
and identify its normalized functional on the selected 032 vectors.
The concrete all-place tower fails the canonical node-gap and
signature requirements. The abstract two-step case is more flexible:
with global Witt index zero, one can take $m=2$ and the Bessel
unipotent rank zero. Its middle space must be totally definite.
A7's natural multiplicity space is totally definite precisely when
the selected first-list CM types agree; that direct realization then
gives a divisor-generated balanced line. B16 supplies a conditional
pole criterion for a different coherent definite source.
The arithmetic kernel/covolume and terminal-mass equations,
and an arithmetic interpretation of the analytic correction,
are the remaining matching problems. B17's upper space is
$A\perp\mathbb H$, of signature $(3,1)$ at every real place;
it is not a hidden definite inner form.

\paragraph{G2: turning the explicit certificates into a uniform
arithmetic lower bound for the prescribed iterated family.}
Q1--Q11 and B5--B7 give actual formulas under specified inputs.
The finite unipotent search yields $|J|\geq1/(4S_{\rm ext})$ with
an explicit grid size after its Gaussian envelope and lattice basis
are supplied. An arithmetic Gram bound needs a coefficient field,
denominator and bounds at all embeddings. Q9 provides resultant
denominators for specified integral Hecke matrices. Q11 transfers
the definite period-square certificate directly to the geometric
ball norm using explicit injective-chart constants; no unknown
ball Gram upper bound is needed. The remaining
step is the normalization and overlap data of the particular
iterated family in G1. Positivity alone cannot supply it, as B4
and B6 show. No universal cycle-degree estimate is asserted.

\paragraph{G3: an exceptional non-CM Hodge summand.}
C16 supplies infinitely many non-CM elliptic source factors, and
C20 adds a geometrically simple genus-two Jacobian with
endomorphism ring $\Z$, on a ball cover of degree at most1764.
C13 gives
an explicit algebraic projector onto a Weil space. One nonzero
projected cycle would generate that two-dimensional rational space.
The known curve-product maps in C12, C16 and C20 are subject to the
same C14 vanishing whenever the target has the generic Weil-type
divisor algebra. The remaining task is a specified map or correspondence
with nonzero Weil projection, or the full-rank pairing data in C04.
C15's characteristic-zero universality implication is a source
reduction, not a proof of that conjecture or of the extraction gate.
The countability result restricts fixed arithmetic towers over
$\C$; it does not preclude the dense algebraic families constructed
here or a further deformation argument.

\paragraph{G4: the chosen integral correspondence and connection.}
C17 identifies the exact horizontal correspondence and local-model
data needed to compare the 032 image cycle with a Kudla--Rapoport
special cycle. C06--C08 prevent replacing those data by an
undifferentiated positivity claim. C11 proves the sharp integral
height-two derivative estimate. C19 removes its artificial
polynomial-coordinate hypothesis and proves integrality of the
CM de Rham normalization coefficient $e_2^*$.
A8 identifies the actual completed 032 tensor with the matched
canonical harmonic jet, including its embedding labels and
factorials. On a matched CM elliptic factor, A8/C19 give the
actual reduced-theta derivative bounds and their prime-to-$p$
level normalization. C21 proves the actual product decomposition
at primes split in $F$ and inert in $E$, and the strong
mixed-derivative bound for canonically matched equal-degree
ratios with integral sections and unit denominator.
For arbitrary integer tensor degree the elliptic bound has
the weaker positive slope, which C22 proves optimal through
exact infinite subsequences on two actual tame theta sections.
Thus transferring the ratio slope to all individual sections
is decided negatively in this normalization.
The general positive-degree higher-dimensional theta estimate
and the particular horizontal Kudla--Rapoport correspondence
remain to be constructed. Neither a raw analytic
frame nor a vertical local intersection is silently used in
that role.

These gates are recorded with dated changes in \texttt{LEDGER.md}.
The mandated three-hour progress review has not been replaced by
this mathematical inventory. Its timing is controlled by the actual
run clock and recorded separately.
